Our results demonstrate that In-PI-MGN, a physics-constrained graph neural network, can accurately predict spatio-temporal blood flow fields across a diverse range of 3D aneurysmal geometries. The model maintains close agreement with CFD ground truth while offering significant computational gains. The model can simulate 80 timesteps in under a minute, whereas traditional CFD methods require approximately one hour for the same task. Even accounting for the 20-hour training time, the proposed approach becomes more efficient after just 20 simulations. \firstround Including dataset generation, the total model development time amounts to approximately 70 hours. \color{black} This leap in speed could be critical in clinical decision support, especially for time-sensitive or resource-constrained contexts.

Indeed, the ability to produce fast full-field hemodynamic predictions opens new possibilities for simulation-augmented clinical workflows. For instance, rupture risk assessments or device deployment planning could benefit from immediate feedback, without the delays of traditional CFD. We foresee this approach contributing to preoperative planning, follow-up monitoring, and even in-room guidance during interventional neuroradiology.

Our contributions are threefold. First, we tackle the challenge of predicting transient, non-Newtonian 3D blood flow, with the ability to accurately capture the evolution of recirculation and swirling patterns, within complex aneurysmal geometries, an area previously unexplored with GNNs. Second, we introduce two architectural advances: contextual inflow feature enhancement to improve long-range temporal representation, and physics-informed training to promote stability and physical plausibility across time. Third, we release BenchAnXplore, a large-scale benchmark dataset of 105 intracranial aneurysm geometries with matched CFD simulations, establishing a new standard for AI-based hemodynamics research.

Prior ML-based approaches for hemodynamic modeling have mostly focused on reduced-order representations, simplified vascular segments, or static features \secondround \cite{pegolotti_learning_2024, li_prediction_2021_1, sen_physics_2024}\color{black}. Our work departs from this by directly learning spatio-temporal dynamics from full 3D velocity fields across realistic geometries. Leveraging the expressive capabilities of GNNs, our model captures complex flow behavior—including swirling and recirculation—within topologically diverse aneurysm shapes. The integration of physics-informed loss terms and contextual inflow features marks a methodological advance, enabling accurate long-term rollout and reducing error accumulation.

\secondround 
Furthermore, we have confronted our method to out-of-distribution unseen cases with different inflow rates and patient-specific cases, and observed that our model had obtained generalization capabilities and was able to uncover realistic flow patterns with no additional training. \color{black} These preliminary results highlight the potential of our approach as a foundation for future fine-tuning toward patient-specific clinical applications.

In summary, this work offers a reproducible, extensible framework that brings high-fidelity hemodynamic inference closer to real-time clinical use. By enabling rapid hemodynamic inference with strong physical consistency, this work paves the way for real-time risk assessment, personalized treatment planning, and simulation-augmented clinical workflows. It offers a reproducible and extensible step toward bridging high-fidelity fluid mechanics and deep learning in cerebrovascular medicine.

\secondround  However, several important limitations remain before clinical deployment can be fully realized. First and foremost, the proposed model is trained on idealized vessel geometries. Clinical routine is characterized by a wide variety of anatomical presentations in terms of size, shape, and orientation. As highlighted by the promising yet degraded performance of our model on out-of-distribution patient-specific cases \thirdround (e.g., average 50-step relative velocity error on the four cases of 11\% compared to 1.6\% on the original test set), \secondround generalization to clinical settings requires training on representative data. While our current dataset is varied and complex enough to serve as a valuable benchmark for model evaluation, it represents only a first step toward patient-specific hemodynamics prediction. In future work, this will need to be complemented by a larger collection of patient-derived geometries or the generation of more realistic synthetic cases.


Second, our CFD ground-truth simulations were generated using a single, temporally and spatially idealized inflow profile. We have demonstrated the model’s ability to generalize across different inflow time profiles, even when these were not part of the training data. Although these inflow rates are widely used in CFD studies of intracranial aneurysms~\cite{cebral_characterization_2005, larrabide_change_2015, lu_hemodynamics_2025}, they remain statistical means extracted from literature~\cite{Reymond2009, Hoi2011}, and may not reflect true patient-specific variability. It would be highly valuable to extend this approach to include clinically acquired Doppler ultrasound or 4D Flow MRI inflow profiles. Additionally, all inlet profiles in this study are prescribed as idealized parabolic flows, which—while common~\cite{qian_risk_2011, larrabide_change_2015, lu_hemodynamics_2025}, omit higher-order effects captured in more realistic Womersley profiles~\cite{cebral_characterization_2005}. Fortunately, previous studies have shown that aneurysmal flow patterns remain robust across such inlet assumptions when appropriate domain lengths are used~\cite{saalfeld_flow-splitting-based_2019}. Nonetheless, we aim to integrate 4D Flow MRI-derived boundary conditions in future work~\cite{brindise_multi-modality_2019}, both for fine-tuning and evaluating our model under more realistic conditions.



Lastly, we acknowledge that while hemodynamic quantities such as wall shear stress (WSS) and oscillatory shear index (OSI) are commonly studied, a biomarker strongly correlated with intracranial aneurysm rupture risk remains elusive. The interpretation of these biomarkers, including proposed thresholds for risk stratification, continues to be debated within the scientific and medical communities. While the primary objective of this work is not direct rupture prediction, it is important to clarify that our model outputs are intended to serve as high-fidelity, physics-consistent inputs to further risk analysis, not as standalone diagnostic tools. We see this as a step toward augmenting clinical decision-making by providing accessible, real-time flow field estimation, which can be combined with other risk factors to support informed diagnostic. \color{black}

Building on this foundation, future work may explore advanced architectures such as multigrid GNNs \cite{garnier_multi-grid_2024}, which enable hierarchical information flow across mesh scales, and attention-based GNNs, which enhance the model’s ability to prioritize key flow structures. These approaches could enable higher-resolution predictions while preserving computational efficiency, further advancing AI-driven hemodynamic modeling. \firstround Moreover, while we demonstrate zero-shot generalization on both unseen inflow profiles and a fully patient-specific case, we emphasize the need for further validation on larger, more heterogeneous clinical datasets, to support further research. \color{black}